\documentclass[12pt]{article}
\usepackage[a4paper,margin=1in]{geometry}
\usepackage[utf8]{inputenc}
\usepackage[T1]{fontenc}
\usepackage{amsmath,amssymb}
\usepackage{textcomp}
\usepackage{parskip}
\begin{document}
\begin{center}
{\Large \textbf{2022 STEP III}}
\end{center}

\subsection*{Page 1}
\begin{flushleft}
\bigskip\textbf{Section A: Pure Mathematics}\par
\bigskip\textbf{Question 1.} Let C 1 be the curve given by the parametric equations\par
\[x = ct , y = ,\]
t\par
\[where c > 0 and t \ne 0, and let C 2 be the circle\]
\[(x - a) 2 + (y - b) 2 = r^2 .\]
C 1 and C 2 intersect at the four points P i (i = 1, 2, 3, 4), and the corresponding values of the\par
parameter t at these points are t i.\par
\textbf{(i)} Show that t i are the roots of the equation\par
\[c2t4 - 2act3 + (a^2 + b^2 - r^2 )t^2 - 2bct + c^2 = 0 . (∗)\]
\textbf{(ii)} Show that\par
\[\leq4\]
\[t^2 = (a^2 - b^2 + r^2 )\]
\[i c^2\]
\[i=1\]
\[\leq4\]
and find a similar expression for .\par
\[t^2\]
\[i=1 i\]
\[\leq4\]
\textbf{(iii)} Hence show that OP i 2 = 4r2 , where OP i denotes the distance of the point P i from\par
\[i=1\]
the origin.\par
\textbf{(iv)} Suppose that the curves C 1 and C 2 touch at two distinct points.\par
By considering the product of the roots of (∗), or otherwise, show that the centre of\par
\[circle C 2 must lie on either the line y = x or y = -x.\]
\end{flushleft}
\subsection*{Page 2}
\begin{flushleft}
\bigskip\textbf{Question 2.} (i) Suppose that there are three non-zero integers a, b and c for which a3 + 2b3 + 4c3 = 0.\par
\[Explain why there must exist an integer p, with |p| < |a|, such that 4p3 + b^3 + 2c3 = 0,\]
\[and show further that there must exist integers p, q and r, with |p| < |a|, |q| < |b| and\]
\[|r| < |c|, such that p^3 + 2q3 + 4r3 = 0. Deduce that no such integers a, b and c can\]
exist.\par
\textbf{(ii)} Prove that there are no non-zero integers a, b and c for which 9a3 + 10b3 + 6c3 = 0.\par
\textbf{(iii)} By considering the expression (3n ± 1) 2 , prove that, unless an integer is a multiple of\par
three, its square is one more than a multiple of 3. Deduce that the sum of the squares\par
of two integers can only be a multiple of three if each of the integers is a multiple of\par
three.\par
\[Hence prove that there are no non-zero integers a, b and c for which a^2 + b^2 = 3c2 .\]
\textbf{(iv)} Prove that there are no non-zero integers a, b and c for which a2 + b2 + c2 = 4abc.\par
\bigskip\textbf{Question 3.} (i) The curve C 1 has equation\par
\[ax2 + bxy + cy2 = 1\]
\[where abc \ne 0 and a > 0.\]
\[Show that, if the curve has two stationary points, then b^2 < 4ac .\]
\textbf{(ii)} The curve C 2 has equation\par
\[ay3 + bx2y + cx = 1\]
\[where abc \ne 0 and b > 0.\]
Show that the x-coordinates of stationary points on this curve satisfy\par
\[4cb3x4 - 8b3x3 - ac3 = 0 .\]
\[Show that, if the curve has two stationary points, then 4ac6 + 27b3 > 0 .\]
\textbf{(iii)} Consider the simultaneous equations\par
\[ay3 + bx2y + cx = 1\]
\[2bxy + c = 0\]
\[3ay2\]
+\par
\[bx2\]
\[= 0\]
\[where abc \ne 0 and b > 0 .\]
Show that, if these simultaneous equations have a solution, then\par
\[4ac6\]
+\par
\[27b3\]
\[= 0 .\]
\end{flushleft}
\subsection*{Page 3}
\begin{flushleft}
\bigskip\textbf{Question 4.} You may assume that all infinite sums and products in this question converge.\par
\textbf{(i)} Prove by induction that for all positive integers n,\par
\[\geq ( \geq ( \geq ( \geq (\]
x x x x\par
\[sinh x = 2 n cosh cosh \cdots cosh sinh\]
2 4 2n 2n\par
and deduce that, for x \ne 0,\par
x\par
\[\geq ( \geq ( \geq (\]
sinh x 2n x x x\par
\[x \geq x ( = cosh 2 cosh 4 \cdots cosh 2n .\]
sinh\par
2n\par
\textbf{(ii)} You are given that the Maclaurin series for sinh x is\par
\[\leq\infty x2r+1\]
\[sinh x = .\]
(2r + 1)!\par
\[r=0\]
y\par
Use this result to show that, as y tends to 0, tends to 1.\par
sinh y\par
Deduce that, for x \ne 0,\par
\[\geq ( \geq ( \geq (\]
sinh x x x x\par
\[= cosh cosh \cdots cosh \cdots .\]
x 2 4 2n\par
\[\geq (\]
x\par
\textbf{(iii)} Let x = ln 2. Evaluate cosh and show that\par
\[\geq x ( 1 + 2 1 2\]
\[cosh = .\]
\[4 2 \times 2 1 4\]
Use part (ii) to show that\par
1 1 1\par
1 1 + 2 2 1 + 2 4 1 + 2 8\par
\[= \times \times \cdots .\]
ln 2 2 2 2\par
\textbf{(iv)} Show that\par
\[\int\]
\[\sqrt ) \sqrt ) \sqrt\]
2 2 2 + 2 2 + 2 + 2\par
\[= \times \times \cdots .\]
\[\pi 2 2 2\]
\end{flushleft}
\subsection*{Page 4}
\begin{flushleft}
\bigskip\textbf{Question 5.} (i) Show that \sum\par
a\par
\[dx = a for all a \leq 0.\]
1 + ex\par
-a\par
\textbf{(ii)} Explain why, if g is a continuous function and\par
\[\sum\]
a\par
\[g(x) dx = 0 for all a \leq 0,\]
\[then g(x) = 0 for all x \leq 0.\]
\[Let f be a continuous function with f(x) \leq 0 for all x. Show that\]
\[\sum\]
a\par
\[dx = a for all a \leq 0\]
1 + f(x)\par
-a\par
if and only if\par
1 1\par
\[+ - 1 = 0 for all x \leq 0,\]
1 + f(x) 1 + f(-x)\par
\[and hence if and only if f(x)f(-x) = 1 for all x.\]
\textbf{(iii)} Let f be a continuous function such that, for all x, f(x) \leq 0 and f(x)f(-x) = 1. Show\par
\[that, if h is a continuous function with h(x) = h(-x) for all x, then\]
\[\sum \sum\]
a h(x) a\par
\[dx = h(x) dx .\]
1 + f(x)\par
-a 0\par
\textbf{(iv)} Hence find the exact value of\par
\[\sum\]
\[1 2 \pi e -x cos x\]
dx .\par
\[- 1 \pi cosh x\]
\end{flushleft}
\subsection*{Page 5}
\begin{flushleft}
\bigskip\textbf{Question 6.} (i) Show that when α is small, cos(θ + α) - cos θ \approx -α sin θ - 1 α2 cos θ.\par
Find the limit as α \to 0 of\par
sin(θ + α) - sin θ\par
(∗)\par
cos(θ + α) - cos θ\par
in the case sin θ \ne 0.\par
\[In the case sin θ = 0, what happens to the value of expression (∗) when α \to 0?\]
\textbf{(ii)} A circle C 1 of radius a rolls without slipping in an anti-clockwise direction on a fixed\par
circle C 2 with centre at the origin O and radius (n - 1)a, where n is an integer greater\par
t(han 2. The) point P is fixed on C 1. Initially the centre of C 1 is at (na, 0) and P is at\par
(n + 1)a, 0 .\par
(a) Let Q be the point of contact of C 1 and C 2 at any time in the rolling motion.\par
Show that when OQ makes an angle θ, measured anticlockwise, with the\par
\[positive x-axis, the x-coordinate of P is x(θ) = a (n cos θ + cos nθ), and find\]
the corresponding expression for the y-coordinate, y(θ), of P .\par
(b) Find the values of θ for which the distance OP is (n - 1)a.\par
\[(c) Let θ 0 = n - 1 \pi. Find the limit as α \to 0 of\]
y(θ 0 + α) - y(θ 0) .\par
x(θ 0 + α) - x(θ 0)\par
( )\par
Hence show that, at the point x(θ 0) , y(θ 0) , the tangent to the curve traced\par
out by P is parallel to OP .\par
\end{flushleft}
\subsection*{Page 6}
\begin{flushleft}
\bigskip\textbf{Question 7.} Let n be a vector of unit length in three dimensions. For each vector r, f(r) is defined by\par
\[f(r) = n \times r .\]
\textbf{(i)} Given that ⎛ ⎞ ⎛ ⎞\par
a x\par
\[n = ⎝ b ⎠ and r = ⎝ y ⎠ ,\]
c z\par
\[show that the x-component of f(f(r)) is -x(b^2 + c^2 ) + aby + acz . Show further that\]
\[f(f(r)) = (n.r)n - r.\]
Explain, by means of a diagram, how f(f(r)) is related to n and r.\par
\textbf{(ii)} Let R be the point with position vector r and P be the point with position vector g(r),\par
where g is defined by\par
\[g(s) = s + sin θ f(s) + (1 - cos θ) f(f(s)) .\]
By considering g(n) and g(r) when r is perpendicular to n, state, with justification,\par
the geometric transformation which maps R onto P .\par
\textbf{(iii)} Let R be the point with position vector r and Q be the point with position vector h(r),\par
where h is defined by\par
\[h(s) = -s - 2 f(f(s)) .\]
State, with justification, the geometric transformation which maps R onto Q.\par
\bigskip\textbf{Question 8.} (i) Use De Moivre's theorem to prove that for any positive integer k > 1,\par
\alpha \alpha \beta \alpha \beta \beta\par
k k\par
\[sin(kθ) = sin θ cos k-1 θ k - (sec 2 θ - 1) + (sec 2 θ - 1) 2 - \cdots\]
3 5\par
and find a similar expression for cos(kθ).\par
\textbf{(ii)} Let θ = cos -1 ( 1 ), where θ is measured in degrees, and a is an odd integer greater\par
a\par
than 1.\par
\[Suppose that there is a positive integer k such that sin(kθ) = 0 and sin(mθ) \ne 0 for\]
\[all integers m with 0 < m < k.\]
\[Show that it would be necessary to have k even and cos( 1 kθ) = 0.\]
Deduce that θ is irrational.\par
\textbf{(iii)} Show that if φ = cot -1 ( 1 ), where φ is measured in degrees, and b is an even integer\par
b\par
greater than 1, then φ is irrational.\par
\end{flushleft}
\subsection*{Page 7}
\begin{flushleft}
\bigskip\textbf{Section B: Mechanics}\par
9 (i) Two particles A and B, of masses m and km respectively, lie at rest on a smooth\par
horizontal surface. The coefficient of restitution between the particles is e, where\par
\[0 < e < 1. Particle A is then projected directly towards particle B with speed u.\]
Let v 1 and v 2 be the velocities of particles A and B, respectively, after the collision, in\par
the direction of the initial velocity of A.\par
1 - ke 1 + e\par
\[Show that v 1 = αu and v 2 = βu, where α = k + 1 and β = k + 1 .\]
Particle B strikes a vertical wall which is perpendicular to its direction of motion and\par
a distance D from the point of collision with A, and rebounds. The coefficient of\par
restitution between particle B and the wall is also e.\par
Show that, if A and B collide for a second time at a point 1 D from the wall, then\par
\[1 + e - e^2\]
\[k = .\]
e(2e + 1)\par
\textbf{(ii)} Three particles A, B and C, of masses m, km and k2m respectively, lie at rest on a\par
smooth horizontal surface in a straight line, with B between A and C. A vertical wall\par
is perpendicular to this line and lies on the side of C away from A and B. The distance\par
between B and C is equal to d and the distance between C and the wall is equal to\par
3d. The coefficient of restitution between each pair of particles, and between particle\par
\[C and the wall, is e, where 0 < e < 1. Particle A is then projected directly towards\]
particle B with speed u.\par
Show that, if all three particles collide simultaneously at a point 3 d from the wall,\par
\[then e = 1 .\]
\bigskip\textbf{Question 10.} Two light elastic springs each have natural length a. One end of each spring is attached to\par
a particle P of weight W . The other ends of the springs are attached to the end-points, B\par
and C, of a fixed horizontal bar BC of length 2a. The moduli of elasticity of the springs PB\par
and PC are s 1 W and s 2 W respectively; these values are such that the particle P hangs in\par
\textasciicircum{}\circ\par
equilibrium with angle BPC equal to 90 .\par
sin θ\par
\textbf{(i)} Let angle PBC = θ. Show that s 1 = 2 cos θ - 1 and find s 2 in terms of θ.\par
\textbf{(ii)} Take the zero level of gravitational potential energy to be the horizontal bar BC and\par
let the total potential energy of the system be -paW . Show that p satisfies\par
\[\sqrt \sqrt\]
1 1\par
\[2 \leq p > (1 + 3)\]
2 4\par
\[and hence that p = 0.7, correct to one significant figure.\]
\end{flushleft}
\subsection*{Page 8}
\begin{flushleft}
\bigskip\textbf{Section C: Probability and Statistics}\par
\bigskip\textbf{Question 11.} A fair coin is tossed N times and the random variable X records the number of heads. The\par
mean deviation, δ, of X is defined by\par
( )\par
\[δ = E |X - μ|\]
where μ is the mean of X.\par
\textbf{(i)} Let N = 2n where n is a positive integer.\par
( )\par
\[(a) Show that P(X \geq n - 1) = 1 1 - P(X = n) .\]
(b) Show that\par
\alpha \beta\par
\[n\leq-1\]
2n 1\par
\[δ = (n - r) .\]
r 22n-1\par
\[r=0\]
\[(c) Show that for r > 0, \alpha \beta \alpha \beta\]
2n 2n - 1\par
\[r = 2n .\]
r r - 1\par
Hence show that \alpha \beta\par
n 2n\par
\[δ = .\]
22n n\par
\textbf{(ii)} Find a similar expression for δ in the case N = 2n + 1.\par
\end{flushleft}
\subsection*{Page 9}
\begin{flushleft}
\bigskip\textbf{Question 12.} (i) The point A lies on the circumference of a circle of radius a and centre O. The\par
point B is chosen at random on the circumference, so that the angle AOB has a\par
\[uniform distribution on [ 0 , 2\pi]. Find the expected length of the chord AB.\]
\textbf{(ii)} The point C is chosen at random in the interior of a circle of radius a and centre O, so\par
that the probability that it lies in any given region is proportional to the area of the\par
region. The random variable R is defined as the distance between C and O.\par
Find the probability density function of R.\par
Obtain a formula in terms of a, R and t for the length of a chord through C that makes\par
an acute angle of t with OC.\par
Show that as C varies (with t fixed), the expected length L(t) of such chords is given\par
by\par
4a(1 - cos 3 t)\par
\[L(t) = .\]
3 sin 2 t\par
Show further that\par
4a ( )\par
\[L(t) = cos t + 1 sec 2 ( 1 t) .\]
3 2 2\par
\textbf{(iii)} The random variable T is uniformly distributed on [ 0 , 1 \pi]. Find the expected value\par
of L(T ).\par
\end{flushleft}
\end{document}
